\begin{prob}[\eligible]
    Suppose you need to keep track of a large dynamic collection of numbers.
    The three most common operations you will be performing are insertions, deletions,
    and queries for the maximum element in the collection. You may assume that the number
    of maximum queries you will be making is close to the number of insertions/deletions.
    You will also be performing membership queries.

    Which data structure would you use:
    a (balanced) binary search tree, or a hash table?  Justify your answer by
    comparing the time complexity of max queries, membership queries, and
    insertions/deletions between the two options.  You may assume that the expected
    time complexity is most important for your application. Your solution should use 
    a reasonable amount of memory and that it should work for negative numbers and decimals.

    \begin{soln}
        We should use a balanced binary search tree.

        While the balanced binary search tree is marginally slower for insertions,
        deletions, and membership queries ($\Theta(\log n)$ vs. $\Theta(1)$ expected time), it is
        \emph{much} faster for maximum queries ($\Theta(\log n)$ vs $\Theta(n)$). 
        Since the number of maximum queries is comparable to the number of insertions/deletions,
        the binary search tree will be faster in aggregate.
    \end{soln}

\end{prob}
